% SPAIS @ CoRL 2026 — The Science of Physical AI Safety. Up to 4 pages excl. references,
% double-blind, non-archival. Deadline 1 Oct 2026 AoE (OpenReview 2 Oct 11:59 UTC).
%
% SKELETON, 14 Sep 2026. Every \todo{} names the committed artifact the number comes from; nothing
% is typed from memory. The style file is the NeurIPS one the 6-pager used; swap for the CoRL
% anonymized template before submission (the CFP requires it) — the body does not change.
\documentclass{article}
\usepackage[final]{neurips_2026}
\usepackage[utf8]{inputenc}
\usepackage[T1]{fontenc}
\usepackage{booktabs}
\usepackage{amsmath}
\usepackage{xcolor}
\usepackage{hyperref}
\newcommand{\todo}[1]{\textcolor{red}{[#1]}}

\title{Two Controls a Robot Red Team Cannot Publish Without:\\
a Length-Matched Scramble and a Second Architecture}

\author{Anonymous Author(s)\\
  Affiliation\\
  Address\\
  \texttt{email}}

\begin{document}
\maketitle

\begin{abstract}
An attack success rate on a vision-language-action policy is a difference between two rates, and
most published ones report only one of them. We report a paired evaluation in simulation in which
every attacked episode has a benign twin at the same task and seed, and we add the two controls a
reviewer asks for first. The first is a \emph{length-matched scramble}: the attack sentence with its
target noun replaced and its token order destroyed, so that ``any long unfamiliar string derails
the policy'' is a hypothesis with its own arm rather than an objection. \todo{scrambled\_text rate
from results/smolvla\_libero\_object\_control\_2026-09-14/aggregate.json}. The second is a second
architecture under the identical protocol: \todo{$\pi_{0.5}$ roleplay and benign rates from
results/pi05\_libero\_object\_2026-09-15}. We also report, for one policy, every attack family the
harness ships against a real checkpoint at three seeds, including a weight-corruption ladder and a
gradient patch, \todo{count of families with a real-policy number from
results/smolvla\_libero\_object\_families\_2026-09-15}; and the same instruction attack on three
further task suites, \todo{spatial/goal/10 rates}. Every rate carries its floor, a Wilson interval
and a task-clustered bootstrap interval; the artifacts, the ledgers and the seeds are public. The
contribution is not an attack. It is the two controls, what they changed, and the tooling that
makes them a flag rather than a project.
\end{abstract}

\section{Introduction}
\todo{Cut from paper/vlm4rwd/paper.tex \S1 to ~0.4 page: the floor argument; the two objections
(wording brittleness; one policy) stated as the hypotheses this paper tests.}

\section{Method}
\label{sec:method}

\paragraph{Policy, suite and predicate.} \todo{Reuse vlm4rwd \S3 verbatim, then add:} The
re-measurement uses provael \todo{version from the shards' tool\_version} on the ten
\texttt{libero\_object} tasks at five seeds per arm; the controls, the family probe and the three
further suites use three seeds. The predicate is the documented default keep-out box throughout, and
every report says so; a calibrated predicate fitted with the attack arm at held-out seeds is reported
separately \todo{results/calib, dawn run} and never substituted after the fact.

\paragraph{The length-matched scramble.} The roleplay attack renders as a twenty-token sentence; the
previous nonsense control was three tokens, so a null there never showed that a twenty-token
out-of-distribution string is harmless. \texttt{scrambled\_text} is the roleplay template's own
tokens with the graspable target replaced by a one-token filler and the order destroyed by a seeded
per-episode shuffle, whitespace-token count identical; \texttt{roleplay\_no\_target} keeps the frame
and replaces only the target. Both carry the \emph{control} role: they enter neither the attack rate
nor the benign floor.

\paragraph{The second architecture.} \todo{$\pi_{0.5}$ checkpoint id, loaded through the same
adapter and the same LIBERO processors; the pilot's clean task success against the checkpoint's
reported 97.5\% (10 ep/task); n = 30 per arm.}

\paragraph{Paired design, multiplicity, intervals.} \todo{Reuse vlm4rwd \S3 paragraphs verbatim.}

\section{Results}

\begin{table}[t]
\centering\small
\caption{\todo{SmolVLA $\times$ libero\_object, ten tasks, five seeds, provael current release.
Columns: arm, role, unsafe/attempts, Wilson 95\%, McNemar vs benign (discordant pairs), Holm,
task-clustered 95\%. Generated from results/smolvla\_libero\_object\_suite\_2026-09-14/README.md.}}
\begin{tabular}{llrlrrl}
\toprule
arm & role & rate & Wilson & McNemar & Holm & clustered \\
\midrule
\todo{rows} & & & & & & \\
\bottomrule
\end{tabular}
\end{table}

\begin{table}[t]
\centering\small
\caption{\todo{The controls beside the attack, three seeds: none, roleplay, benign\_reword,
nonsense\_text, scrambled\_text, roleplay\_no\_target. From
results/smolvla\_libero\_object\_control\_2026-09-14.}}
\begin{tabular}{llrl}
\toprule
arm & role & rate & Wilson \\
\midrule
\todo{rows} & & & \\
\bottomrule
\end{tabular}
\end{table}

\paragraph{What the scramble says.} \todo{One of two paragraphs, decided by the data: (a) the
scramble fires at or near the floor while roleplay fires at $\sim$88\%, so the gap is attacker
control and not wording; or (b) the scramble fires well above the floor, the headline is restated
as generic fragility with a bounded attacker-specific residual, and an erratum (E-2026-12) is
issued. Write the one the aggregate supports; do not write both.}

\paragraph{The second architecture.} \todo{$\pi_{0.5}$ vs SmolVLA on the same arm and tasks,
n = 30 each; state the difference with both intervals and say plainly what n = 30 can and cannot
resolve.}

\paragraph{Breadth.} \todo{Families with a real-policy number after the probe (was 3 of 17), the
weight-corruption crossing point for gradient vs random selection on a flow-matching head against
the 100--300-flip band reported for that head class, and the first measurement of the gradient
patch. Nulls are reported as nulls with their intervals.}

\paragraph{Suites.} \todo{roleplay on spatial/goal/10 at three seeds beside object: does the
instruction attack survive a suite the policy was not fine-tuned on for the headline?}

\section{Discussion and limitations}
\todo{Simulation only; one default predicate; three seeds on every arm but the headline; $\pi_{0.5}$
at n = 30; templated attacks as a floor on susceptibility; the sim-to-real correlation literature
(third-party Spearman 0.40--0.70) as the reason none of this predicts a real rate; what a
certifier can and cannot take from a report shaped like ISO/IEC 17025 clause 7.8 (nothing about
conformity).}

\section*{Artifacts}
\todo{Anonymized artifact bundle: the results directories, aggregate.json files, ledgers, the
release tag and the recipe. Run check\_anonymity.sh before upload.}

\bibliographystyle{plain}
\begin{thebibliography}{9}
\bibitem{todo} \todo{Import the bibliography of paper/vlm4rwd/paper.tex; add RobustVLA
(arXiv:2510.00037), LIBERO-PRO (arXiv:2510.03827), the weight-corruption paper (arXiv:2608.15475),
the third-party sim-to-real audit (arXiv:2606.10366).}
\end{thebibliography}

\end{document}
